% ====================================================================
%  Chapter 2 — Related Work
% --------------------------------------------------------------------
%  Template chapter. The section headings cover the four sub-fields
%  the thesis touches:
%    2.1 Speaker verification architectures
%    2.2 Loss functions for speaker discrimination
%    2.3 Cross-lingual / low-resource transfer
%    2.4 Score normalisation and calibration
%    2.5 Multilingual self-supervised speech representation
%    2.6 Positioning of this work
%  Citations resolve against the .bib block in RUN_GUIDE.md §16.
%  Add references to .bib as you cite them; bracketed placeholders
%  like [CITE_DESPLANQUES_PAPER_PAGE] mark spots where a specific page
%  number or quoted figure should be added during the lit-review pass.
% ====================================================================

\chapter{Related Work}
\label{ch:related}

\section{Speaker verification architectures}
\label{sec:related-architectures}

The dominant architectures for modern speaker verification fall into
three families. The first is the residual convolutional family,
typified by ResNet-SE34V2 with squeeze-and-excitation blocks and
self-attentive pooling. The second is the time-delay neural network
(TDNN) family, dominated since 2020 by ECAPA-TDNN
\cite{desplanques2020ecapa}. ECAPA-TDNN introduced three key
ingredients: (i) Res2Net-style multi-scale convolutions with
hierarchical accumulation, (ii) squeeze-and-excite channel
re-weighting at each block, and (iii) multi-layer feature
aggregation followed by channel-attentive statistical pooling.
ECAPA-TDNN consistently reports 0.8--1.5\,\% EER on VoxCeleb1-O and
has become the de-facto reference architecture in subsequent SV
papers. The third family is self-supervised, replacing the
mel-front-end with a frozen pretrained transformer encoder such as
WavLM \cite{chen2022wavlm}, XLS-R, or mHuBERT-147; these models
typically outperform mel-based architectures on low-resource
language data by 30--50\,\% relative EER, at the cost of significant
parameter count (95M--300M).

% [CITE_OR_EXPAND: brief mention of x-vector / i-vector predecessors
% to ground the field's evolution; one paragraph is enough for an
% MPhil chapter.]

\section{Loss functions for speaker discrimination}
\label{sec:related-loss}

Speaker-verification training has progressively replaced cross-entropy
softmax with margin-based variants. AM-Softmax and Additive Angular
Margin Softmax (AAM-Softmax) \cite{deng2019arcface} reshape the
classification objective to produce more discriminative embeddings
on the hypersphere by adding a margin to the cosine similarity
between the embedding and the target-class weight vector. The
margin-based losses became standard in SV largely because they
produce embeddings that score well under cosine distance at eval
time — a closer match to how SV systems are actually used than the
intra-class compactness produced by plain softmax.

% [CITE_OR_EXPAND: mention sub-centre AAM-Softmax (\S 3.2 \#10 of the
% analysis doc) as a recent refinement; one paragraph if the SC-AAM
% ablation is included, omit otherwise.]

Auxiliary use of language information has also proven productive,
either during training via multi-task language-ID heads, or at
decision time via language-aware scoring and calibration
\cite{thienpondt2022calibration}.
Adversarial counterparts — Domain-Adversarial Neural Networks
\cite{ganin2015dann} attached to a language or channel classifier
through a gradient-reversal layer — pursue the opposite goal: forcing
the embedding to be \emph{uninformative} about the domain. Both
approaches appear in the SV literature; their relative effectiveness
depends on whether language is a useful signal (multi-task wins)
or a confound (adversarial wins).

\section{Cross-lingual and low-resource transfer}
\label{sec:related-crosslingual}

When a target-language SV corpus is small ($\leq$~10k utterances),
training a deep architecture from scratch overfits catastrophically.
The standard remedy is cross-lingual fine-tuning from a checkpoint
pretrained on a high-resource corpus (typically VoxCeleb1 or
VoxCeleb2). Three knobs control how aggressively the target
language overwrites the pretrained representation:

\begin{enumerate}
    \item \textbf{Selective freezing.} Freezing the front-end
    (SincConv or mel-filterbank parameters) preserves the
    acoustic-feature extractor learned on the source language. For
    very small corpora, freezing the entire encoder and training
    only the projection head is the safest option, and already
    captures most of the self-supervised representation's value
    \cite{chen2022wavlm,sang2025unipet}.
    \item \textbf{Learning-rate scaling.} A standard
    multiplier of 0.1$\times$ to 0.01$\times$ the from-scratch LR
    avoids destabilising the prior.
    \item \textbf{Layer-wise learning-rate decay (LLRD).} A decay
    factor $\gamma \in [0.8, 0.95]$ applied per layer from the
    encoder top to the input front-end, so that deeper layers carry
    transferable acoustic priors while shallower (closer-to-output)
    layers adapt more aggressively. LLRD is universal in the BERT
    and WavLM fine-tuning literature; its application to SV
    fine-tuning is less systematically reported, motivating one of
    this thesis's ablations.
\end{enumerate}

% [CITE_OR_EXPAND: position prior cross-lingual SV papers — e.g. NIST
% SRE multilingual evaluations, MLS, MUSAN-based low-resource SV
% work — relative to the SL setting.]

\section{Score normalisation and calibration}
\label{sec:related-scoring}

Two eval-time mechanisms address the gap between training-time and
deployment-time score distributions. Adaptive Symmetric Score
Normalisation (AS-Norm) \cite{matejka2017analysis} computes
per-trial-side cohort statistics from a fixed set of impostor
embeddings and normalises trial scores against those statistics.
It is universal in the NIST SRE pipeline and typically yields
5--20\,\% relative MinDCF reduction at zero training cost.

Probabilistic Linear Discriminant Analysis (PLDA) \cite{kenny2010,
sizov2014unifying} models the embedding distribution generatively
and replaces cosine scoring with a likelihood ratio test under
``same speaker'' vs ``different speaker'' hypotheses. PLDA was the
dominant SV backend in the i-vector era and remains useful for
short-utterance and channel-mismatched trials, although on modern
discriminatively-trained embeddings the marginal gain over cosine is
smaller. The two-covariance simplification of PLDA
\cite{sizov2014unifying} provides nearly identical EER to the full
formulation at half the computational complexity.

Score normalisation and PLDA compose: AS-Norm normalises against
cohort statistics regardless of whether the underlying score is
cosine or PLDA. The NIST SRE pipeline order is PLDA $\to$ AS-Norm.

% [CITE_OR_EXPAND: per-language threshold calibration is a separate
% research direction (Platt scaling, isotonic regression). One short
% paragraph for completeness.]

\section{Multilingual self-supervised speech representation}
\label{sec:related-ssl}

Self-supervised models trained on large multilingual audio corpora
(wav2vec 2.0 \cite{baevski2020wav2vec2}, WavLM
\cite{chen2022wavlm}, XLS-R, mHuBERT-147 \cite{boito2024mhubert})
provide speech-aware embeddings without language-specific labelling.
mHuBERT-147 is particularly relevant to the Sri Lankan setting: it
deliberately upweights low-resource languages including Sinhala and
Tamil in its pretraining corpus.

There are three ways to use such models for low-resource SV:
(i) load an off-the-shelf checkpoint, freeze it, and train only a
pooling head and projection on the target-language labelled data
% [CITE_OR_EXPAND: typical SV-on-WavLM results]; (ii) continue
pretraining (CPT) on unlabelled target-language audio before
SV-fine-tuning, which adds 5--15\,\% relative EER on top of (i)
when sufficient unlabelled audio ($\geq$~500\,h) is available
\cite{pratap2024mms}; (iii) pretrain from scratch on
target-language-only audio, which requires thousands of hours and
rarely outperforms (ii). This thesis evaluates option (i) as the
SSL baseline; option (ii) is documented as future work.

\section{Positioning of this work}
\label{sec:positioning}

The contributions outlined in
Section~\ref{sec:contributions} sit at the intersection of three
otherwise separately-developed lines of work:

\begin{itemize}
    \item \textbf{Architectures and losses} for English-scale SV
    (\S~\ref{sec:related-architectures}, \S~\ref{sec:related-loss})
    are applied to low-resource Sinhala/Tamil with no targeted
    benchmark in the literature.
    \item \textbf{Cross-lingual transfer techniques}
    (\S~\ref{sec:related-crosslingual}) are widely studied for ASR
    and translation but less so for SV; their interaction with
    SV-specific losses (AAM-Softmax) and SV-specific eval backends
    (PLDA, AS-Norm) has not been systematically ablated.
    \item \textbf{Multilingual SSL encoders}
    (\S~\ref{sec:related-ssl}) cover Sinhala and Tamil only as part
    of broader low-resource batches, with no published SV evaluation
    on either language.
\end{itemize}

To the best of our knowledge, this is the first study to provide a
reproducible SV benchmark on Sinhala and Tamil with full ablation
isolation of the cross-lingual transfer levers. The benchmark, the
corpus preparation pipeline, and the open-source toolkit
(\S~\ref{sec:contributions}) are released to support future work in
this direction.
